\section{Cái bọc quanh mỗi tầng}
\label{sec:ch07-quanh-tang}

\index{residual!bọc quanh mỗi tầng con}%
Ba mục trước dựng xong phần attention. Phần còn lại của kiến trúc là những thứ
không ai đặt tên bài báo theo: một phép cộng, một phép chuẩn hóa, và một mạng
hai tầng. Chúng đến từ ba dòng nghiên cứu khác nhau và bài lắp cả ba vào mà
không tranh luận gì.

Mỗi tầng con, dù là attention hay không, nằm trong cùng một cái bọc. Mục 3.1
viết ra bằng một dòng:

\begin{equation}
  \mathrm{output} = \layernorm(\vect{x} + \mathrm{Sublayer}(\vect{x})).
  \label{eq:ch07-post-ln}
\end{equation}

Cộng đầu vào vào đầu ra, rồi chuẩn hóa. Hai phép, và mỗi phép có một bài báo
riêng đứng sau.

\begin{bridge}{He và cộng sự 2015: cộng thẳng đầu vào vào đầu ra}
\index{residual!He 2015}%
Bài này ra đời cho mạng tích chập rất sâu, không liên quan gì tới
chuỗi~\autocite{he2016deep}. Vấn đề họ gặp nằm gọn trong hình 1 của bài: trên
CIFAR-10, mạng 56 tầng có lỗi \emph{huấn luyện} cao hơn mạng 20 tầng.

Chỗ đáng dừng là chữ huấn luyện. Overfit không giải thích được kết quả ấy, vì
overfit đẩy lỗi kiểm tra lên trong khi lỗi huấn luyện đi xuống, còn ở đây cả
hai cùng lên. Mục 1 của bài nói thẳng ra:
\enquote{such degradation is not caused by overfitting, and adding more layers
to a suitably deep model leads to higher training error}. Nói cách khác đây
không phải chuyện mô hình học thuộc, mà là chuyện mô hình khó tối ưu.

Cách chữa là đổi cái mỗi khối phải học. Thay vì bắt khối học ánh xạ
\(H(\vect{x})\), cho nó học \(F(\vect{x}) = H(\vect{x}) - \vect{x}\) rồi cộng
\(\vect{x}\) lại ở cuối, tức \tn{kết nối tắt}{residual connection}. Nếu ánh xạ
tốt nhất cho khối ấy là không làm gì cả, khối chỉ cần đẩy \(F\) về 0, và đẩy về
0 là chuyện dễ.

Chỗ liên quan tới sách này là đạo hàm. Đạo hàm của \(\vect{x} + F(\vect{x})\)
theo \(\vect{x}\) là \(\mat{I} + \jac{F}{\vect{x}}\), nên gradient có một đường
đi thẳng qua ma trận đơn vị, không nhân với gì cả. Đó đúng là hình dạng
mục~\ref{sec:ch06-duong-gradient} tìm thấy ở \eqref{eq:ch06-duong-moi}, và đúng
hình dạng của vòng quay lỗi ở chương~\ref{ch:lstm}. Ba kiến trúc khác hẳn nhau,
cùng một mẹo: cho gradient một lối đi mà không có tích Jacobi nào trên đó.
\end{bridge}

\begin{bridge}{Ba, Kiros và Hinton 2016: chuẩn hóa theo từng mẫu}
\index{layer normalization!Ba 2016}%
\tn{Chuẩn hóa theo tầng}{layer normalization} lấy trung bình và phương sai trên
\emph{các đơn vị của một tầng, trong một mẫu huấn luyện}, rồi đưa về trung bình
0 phương sai 1 và nhân lại bằng hai tham số học được~\autocite{ba2016layer}.
Tóm tắt của bài phát biểu nó như một phép chuyển vị:
\enquote{we transpose batch normalization into layer normalization by computing
the mean and variance used for normalization from all of the summed inputs to
the neurons in a layer on a single training case}.

Chỗ khác với batch normalization nằm ở chỗ lấy thống kê, và tóm tắt nêu ba hệ
quả. Thứ nhất, batch norm lấy trên cả một mini-batch nên tác dụng của nó phụ
thuộc vào kích thước batch (\enquote{the effect of batch normalization is
dependent on the mini-batch size}). Thứ hai, layer norm tính y hệt nhau ở lúc
huấn luyện và lúc chạy thật, còn batch norm thì không:
\enquote{Unlike batch normalization, layer normalization performs exactly the
same computation at training and test times}. Thứ ba, với mạng hồi quy thì bài
nói về batch norm rằng \enquote{it is not obvious how to apply it}, còn layer
norm áp vào thẳng được bằng cách tính thống kê riêng ở từng
\tn{bước thời gian}{timestep}.

Hệ quả thứ ba là hệ quả đáng nhớ ở đây, dù bài viết nó cho mạng hồi quy chứ
không cho kiến trúc này: layer norm không nhìn sang mẫu khác, nên độ dài câu và
cách xếp batch không đi vào phép chuẩn hóa được.

Tôi chỉ đọc bài này ở mức tóm tắt. Một chi tiết thư mục đáng ghi: nó có đúng
một bản arXiv, ngày 21 tháng 7 năm 2016, không số DOI và không dòng nào ghi hội
nghị hay tạp chí. Một bài chưa từng qua bình duyệt chính thức, nằm trong gần
như mọi Transformer từng chạy.
\end{bridge}

\subsection*{Thứ tự của hai phép, và vì sao nó thành một chương khác}

\index{layer normalization!post-LN là thứ tự của bài}%
\eqref{eq:ch07-post-ln} chuẩn hóa \emph{sau} khi cộng. Cách gọi bây giờ là
post-LN, và cách gọi ấy tồn tại vì có một cách khác.

Đặt layer norm vào trong, tức \(\vect{x} + \mathrm{Sublayer}(\layernorm(\vect{x}))\),
thì đường tắt của He đi thẳng từ đầu vào tới đầu ra của cả chồng tầng mà không
gặp phép chuẩn hóa nào. Bài viết post-LN, và mọi thứ trong sách này viết theo
bài. Nhưng thứ tự ấy là lý do mục 5.3 của bài phải có một lịch warmup, và
chương~\ref{ch:chi-phi-va-song-song} là chỗ đo hai thứ tự cạnh nhau thay vì
tuyên bố cái nào đúng.

Tôi nói trước ở đây vì \eqref{eq:ch07-post-ln} là loại phương trình dễ chép lại
mà không để ý thứ tự, và thứ tự là toàn bộ nội dung của nó.

\subsection*{Mạng truyền thẳng, và chữ \enquote{theo vị trí}}

\index{FFN!chạy riêng ở từng vị trí}%
Tầng con thứ hai của mỗi tầng encoder là phương trình (2):

\begin{equation}
  \ffn(\vect{x}) = \max(0, \vect{x}\mat{W}_1 + \vect{b}_1)\mat{W}_2 + \vect{b}_2.
  \label{eq:ch07-ffn}
\end{equation}

Một mạng hai tầng với hàm rectified linear unit, viết tắt là ReLU, ở giữa. Bài
đặt \(\dff = 2048\) so với
\(\dmodel = 512\), tức rộng gấp bốn ở giữa rồi thu về.

Chữ đáng đọc là \enquote{position-wise}. Cùng một bộ trọng số chạy ở từng vị
trí, riêng rẽ, không có gì đi ngang. Bài mô tả nó theo cách khác cho cùng một
ý: hai phép tích chập kernel bằng 1.

Đó không phải một chi tiết cài đặt. Nó là chỗ khóa lại một phát biểu về cả kiến
trúc: \emph{trong Transformer, attention là thứ duy nhất chuyển thông tin giữa
các vị trí.} Mọi thứ khác chạy tại chỗ. Tôi viết một test khẳng định đúng câu
ấy, đổi đầu vào ở một vị trí rồi đòi đúng vị trí ấy đổi và bốn vị trí kia đứng
im, vì nếu \eqref{eq:ch07-ffn} lỡ trộn được các vị trí thì mọi thứ
mục~\ref{sec:ch07-vi-tri} nói về vị trí đều mất chỗ dựa.

\index{residual!vì sao mọi tầng con cùng bề rộng}%
Một hệ quả nhỏ nhưng ràng buộc mọi thứ: vì \eqref{eq:ch07-post-ln} cộng
\(\vect{x}\) vào đầu ra của tầng con, hai vế phải cùng bề rộng. Nên mọi tầng
con và cả embedding đều rộng \(\dmodel\), và bài nói thẳng ra thế. \(\dff\)
được phép khác vì nó nằm gọn bên trong \eqref{eq:ch07-ffn} và không chạm vào
phép cộng.
